

\section{Städtetypen Weltweit}
\begin{myexercise}[Gruppenpuzzle zu Stadt-Typen]
Setzen Sie sich in 3er-Gruppen zusammen. Suchen Sie sich jeweils einen der Texte auf den folgenden Seiten aus. Setzen Sie sich in Gruppen zusammen und zeichnen Sie für jeden der Texte einen Stadtplan auf, den Sie sich gegenseitig erklären. Was zeichnet die jeweiligen Städte aus? Wie sind sie räumlich aufgebaut?
\end{myexercise}

\subsection{Die Europäische Stadt}
\begin{quote}
Die Kleinkammerung des europäischen Raumes, seine hohe Besiedlungsdichte und die historische Vielschichtigkeit haben eine solche Formenvielfalt hervorgebracht, die es schwierig macht, gemeinsame Merkmale in einem Modell zusammenzufassen.
Im hier vorgestellten Modell der europäischen Stadt sind die Grunddaseinsfunktionen und die Gebäudetypen als wichtigste Merkmale dargestellt. Die Grundrissstruktur, insbesondere das Stras­sennetz, reicht vielerorts bis in die mittelalterliche Gründungsphase zurück. Die Gebäude zeigen die historische Schichtung mindestens seit der frühen Neuzeit. Und in den Gebäudetypen sind weitgehend auch die Funktionsräume abgebildet.

Wenn wir eine Stadt erstmals besuchen, orientieren wir uns meistens an der Altstadt, die vielerorts seit dem 19. Jahrhundert das Geschäftszentrum bildet, sowie am Bahnhof. Nach aussen folgt die konventionelle Stadtkernbebauung mit sehr hoher Dichte aus der Zeit der Industrialisierung im 19. Jahrhundert. Diese Quartiere entstanden nach der Eröffnung der Eisenbahnlinien in unmittelbarer Umgebung des Bahnhofs. Die neue Bahnhofstrasse entwickelte sich zu einer wichtigen Geschäftsstrasse. Für die rasch wachsende Bevölkerung entstanden die ebenfalls sehr dichten Quartiere mit Blockrandbebauungen. Die 3- bis 5-geschossigen, zusammengebauten Häuser stehen dicht am Trottoir, weisen jedoch einen mehr oder weniger grossen Innenhof auf.

\begin{figure}[H]
\centering
\includegraphics[width=\linewidth]{"Private/Images/EuropaeischeStadtModell"}
\caption{Modell der Europäischen Stadt}
\end{figure}

Entlang der Ausfallstrassen, die zusammen mit den Eisenbahnlinien und den Ringstrassen der neuesten Zeit die wichtigsten Linienelemente bilden, entstanden bereits im 19. Jahrhundert einzelne Wohn- und Gewerbegebäude. Ebenfalls im 19. Jahrhundert entstanden in den bevorzugten Lagen Villenquartiere, als Einfamilien- oder als Doppelhäuser. Ihre Gärten weisen heute oft mächtige Bäume auf, die Quartierstrassen sind nur schmal, da sie noch nicht für den Autoverkehr gebaut wurden. Viele dieser alten Villen werden heute von Dienstleistungsbetrieben genutzt (Anwaltskanzleien, Arztpraxen, Architekturbüros, Beratungsunternehmen).

Nach dem Ersten Weltkrieg wurden am Stadtrand ganze Quartiere mit Reiheneinfamilienhäusern oder mit Mehrfamilienhäusern gebaut. Beide Quartiertypen weisen eine niedrige bis mittlere Dichte auf, viele entstanden als Genossenschaftssiedlungen.

Und schliesslich bilden die Industriequartiere entlang der Eisenbahnlinien typische Elemente vieler europäischer Städte. Diese Areale sind heute wichtige Flächen für die Stadterneuerung. Ab 1950 wurden als Folge der starken Bevölkerungszunahme zahlreiche Wohnquartiere in differenzierter Bauweise mit Hoch- und Scheibenhäusern erbaut. Ab etwa 1965 wurden an südexponierten Hanglagen Terrassensiedlungen erbaut, die eine Mittelstellung zwischen Reiheneinfamilienhaus und mehrgeschossigem Haus einnehmen und eine recht hohe Bebauungsdichte aufweisen.

Für die Lebensqualität besonders wichtig sind die innerstädtischen Grünanlagen (Stadtparks, Friedhofsanlagen) und die nahen Stadtwälder, ohne die eine Stadt nicht existieren und sich entwickeln kann.

-- \textit{Kurz-Beschreibung der europäischen Stadt nach \textcite{hep}}
\end{quote}

Schauen Sie sich auch nochmals die Skripte der Doppellektionen 1 und 2 an, um Ihr Verständnis europäischer Städte aufzufrischen!

\subsection{Die Nordamerikanische Stadt}

\begin{figure}[H]
\centering
\includegraphics[trim={1cm 1.5cm 1cm 2cm},clip,width=\linewidth, page=5]{"Private/Appendix/stadtmodelleDiercke"}
\caption{Kurz-Beschreibung der nordamerikanischen Stadt \cite[354]{westermann}}
\end{figure}
\begin{figure}[H]
\centering
\includegraphics[trim={1cm 1.5cm 1cm 2cm},clip,width=\linewidth, page=6]{"Private/Appendix/stadtmodelleDiercke"}
\caption{Kurz-Beschreibung der nordamerikanischen Stadt \cite[355]{westermann}}
\end{figure}

\subsection{Die Orientalische Stadt}

\begin{figure}[H]
\centering
\includegraphics[trim={1cm 1.5cm 1cm 2cm},clip,width=\linewidth, page=1]{"Private/Appendix/stadtmodelleDiercke"}
\caption{Kurz-Beschreibung der orientalischen Stadt \cite[350]{westermann}}
\end{figure}
\begin{figure}[H]
\centering
\includegraphics[trim={1cm 3cm 1cm 2cm},clip,width=\linewidth, page=2]{"Private/Appendix/stadtmodelleDiercke"}
\caption{Kurz-Beschreibung der orientalischen Stadt \cite[351]{westermann}}
\end{figure}

\begin{figure}[H]	
\centering
\includegraphics[width=.96\linewidth]{"Private/Images/Orient-Stadt/fes"}
\includegraphics[width=.48\linewidth]{"Private/Images/Orient-Stadt/koranschule-fes"}
\includegraphics[width=.48\linewidth]{"Private/Images/Orient-Stadt/riad-rabat"}
\includegraphics[width=.48\linewidth]{"Private/Images/Orient-Stadt/stadttor-fes-2"}
\includegraphics[width=.48\linewidth]{"Private/Images/Orient-Stadt/strasse-rabat"}
\caption{Bilder aus marokkanischen Städten (Fès, Rabat)}
\end{figure}

\subsection{Die Lateinamerikanische Stadt}

\begin{figure}[H]
\centering
\includegraphics[trim={1cm 1.5cm 1cm 2cm},clip,width=\linewidth, page=3]{"Private/Appendix/stadtmodelleDiercke"}
\caption{Kurz-Beschreibung der lateinamerikanischen Stadt \cite[352]{westermann}}
\end{figure}
\begin{figure}[H]
\centering
\includegraphics[trim={1cm 3cm 1cm 2cm},clip,width=\linewidth, page=4]{"Private/Appendix/stadtmodelleDiercke"}
\caption{Kurz-Beschreibung der lateinamerikanischen Stadt \cite[353]{westermann}}
\end{figure}

